\section{Probabilidad del modelo, confianza y calibración}

Las preguntas y respuestas en clase señalan que el promedio de log probability de los tokens de salida se correlaciona con la corrección, pero por lo general la calibración es insuficiente: una respuesta incorrecta también puede recibir una probabilidad interna muy alta, y el modelo sigue defendiendo su respuesta original incluso cuando se le corrige.

Se puede medir con un reliability diagram o con el Expected Calibration Error:

$$
\operatorname{ECE}=\sum_{b=1}^{B}\frac{|S_b|}{N}
\left|\operatorname{acc}(S_b)-\operatorname{conf}(S_b)\right|.
$$

\begin{itemize}
    \item $S_b$: las muestras cuya confianza de predicción cae en el $b$-ésimo intervalo;
    \item $\operatorname{acc}(S_b)$: la tasa de aciertos real de ese intervalo;
    \item $\operatorname{conf}(S_b)$: la confianza promedio de ese intervalo;
    \item Cuanto menor sea el ECE, más cerca está la confianza de la tasa de aciertos empírica.
\end{itemize}

\begin{knowledgebox}{Lo ideal es que el selector aprenda señales verificables, no que solo lea la confianza propia del modelo}
En tareas de código se pueden incorporar compilación, casos de prueba, mutation testing, análisis de complejidad y clustering de comportamiento; en un agent de investigación se pueden incorporar la calidad de la fuente, la consistencia de las citas y la evidencia cruzada. Las señales externas suelen ser más confiables que la confianza autorreportada por el modelo generador.
\end{knowledgebox}

\subsection{Resumen de la sección}

La probabilidad de un modelo de lenguaje no es una «tasa de corrección de la respuesta» calibrada. El ordenamiento de candidatos, las condiciones de paro y si se activa o no la búsqueda deben combinarse con verificación externa y calibración especializada, y no basarse únicamente en la probabilidad cruda de los tokens.
